% --- Archon physics formalization source begin ---
% archon:physics
% archon:covers PhyXMiniProblems/problem_phyx_mini_0144.lean
% archon:source-report reports/phyx_mini/problem_phyx_mini_0144.source.json

\chapter{Physics problem phyx\_mini\_0144}
\label{ch:PhyXMiniProblems_problem_phyx_mini_0144}

\paragraph{Problem source.}
\#\# Physical scenario

A lighted candle is placed \textbackslash{}(36\textbackslash{},\textbackslash{}mathrm\{cm\}\textbackslash{}) in front of a converging lens of focal length \textbackslash{}(f\_1 = 13\textbackslash{},\textbackslash{}mathrm\{cm\}\textbackslash{}), which in turn is \textbackslash{}(56\textbackslash{},\textbackslash{}mathrm\{cm\}\textbackslash{}) in front of another converging lens of focal length \textbackslash{}(f\_2 = 16\textbackslash{},\textbackslash{}mathrm\{cm\}\textbackslash{}) (see figure).

\#\# Question

Calculate the relative size of the final image.

\#\# Answer choices

- A: A: \$0.56 \textbackslash{}times x\$

- B: B: \$0.36 \textbackslash{}times x\$

- C: C: \$0.46 \textbackslash{}times x\$

- D: D: \$0.40 \textbackslash{}times x\$

\#\# Figure caption (auxiliary; use the image as primary evidence)

The image shows a diagram related to optics. It features the following elements:

1. **Candle**: On the left side, there is an illustration of a lit candle with a glow around the flame.

2. **Lenses**: Two convex lenses are displayed:
   - The first lens (on the left) has a focal length labeled as \textbackslash{}( f\_1 = 13 \textbackslash{}, \textbackslash{}text\{cm\} \textbackslash{}).
   - The second lens (on the right) has a focal length labeled as \textbackslash{}( f\_2 = 16 \textbackslash{}, \textbackslash{}text\{cm\} \textbackslash{}).

3. **Distances**: There are horizontal lines with measurements indicating distances:
   - A line labeled "36 cm" shows the distance from the candle to the first lens.
   - A line labeled "56 cm" shows the distance between the first lens and the second lens.

The arrangement indicates a system for a physics or optics problem involving lens positioning relative to a light source.

\#\# Dataset metadata

Geometrical Optics; Multi-Formula Reasoning, Spatial Relation Reasoning

\paragraph{Recorded answer/context.}
C: C: \$0.46 \textbackslash{}times x\$

\paragraph{Figure/image path.}
/root/proposal\_for\_physic/science-mango/phyx\_mini\_run/phyx\_data/test\_image/144.png

\paragraph{Formalization target.}
create a compiling Lean file with sorry bodies at `PhyXMiniProblems/problem\_phyx\_mini\_0144.lean`. The Lean declarations must preserve the physical quantities, dimensions or dimensional roles, figure labels, governing-law hypotheses, and final relation expressed by this problem.
Use Mathlib/Physlib names found through LeanExplore where available. If a physics API is missing, introduce faithful local abstractions rather than scalar placeholder aliases.

\begin{theorem}[Physics formalization target]
\label{thm:physics:phyx_mini_0144:target}
\lean{PhyXMiniProblems.ProblemPhyXMini0144.problem_phyx_mini_0144}
\uses{def:physics:phyx-mini-0144:phyxminiproblems-problemphyxmini0144-twolenscandlesetup, def:physics:phyx-mini-0144:phyxminiproblems-problemphyxmini0144-hasdepictedlensarrangement, def:physics:phyx-mini-0144:phyxminiproblems-problemphyxmini0144-hasstatedfigurereadouts, def:physics:phyx-mini-0144:phyxminiproblems-problemphyxmini0144-hassequentialimagegeometry, def:physics:phyx-mini-0144:phyxminiproblems-problemphyxmini0144-hasphysicalrealimageconfiguration, def:physics:phyx-mini-0144:phyxminiproblems-problemphyxmini0144-satisfiesthinlensequationat, def:physics:phyx-mini-0144:phyxminiproblems-problemphyxmini0144-satisfiestransversesizelawat, def:physics:phyx-mini-0144:phyxminiproblems-problemphyxmini0144-finalrelativesize, def:physics:phyx-mini-0144:phyxminiproblems-problemphyxmini0144-matchesdisplayedrelativesize, def:physics:phyx-mini-0144:phyxminiproblems-problemphyxmini0144-isnearestdisplayedchoice}
The assigned autoformalize agent should translate this physics problem into Lean declarations in the covered file, with theorem and lemma proof bodies written as `by sorry`.
\end{theorem}
\begin{proof}
This is an autoformalization task, not a proof task. Produce statements that can later be proved by the physics prover without weakening the physical model.
\end{proof}

% --- Archon declaration topology begin ---
\section{Declaration topology}
\begin{definition}[Length Magnitude]
  \label{def:physics:phyx-mini-0144:phyxminiproblems-problemphyxmini0144-lengthmagnitude}
  \lean{PhyXMiniProblems.ProblemPhyXMini0144.LengthMagnitude}
  A nonnegative physical length, independent of the unit used to read it
\end{definition}
\begin{proof}
  This is the declared model definition of Length Magnitude.
\end{proof}
\begin{definition}[centimeter Unit Choices]
  \label{def:physics:phyx-mini-0144:phyxminiproblems-problemphyxmini0144-centimeterunitchoices}
  \lean{PhyXMiniProblems.ProblemPhyXMini0144.centimeterUnitChoices}
  Unit choices in which physical lengths are read in centimeters
\end{definition}
\begin{proof}
  This is the declared model definition of centimeter Unit Choices.
\end{proof}
\begin{definition}[centimeters Value]
  \label{def:physics:phyx-mini-0144:phyxminiproblems-problemphyxmini0144-centimetersvalue}
  \lean{PhyXMiniProblems.ProblemPhyXMini0144.centimetersValue}
  \uses{def:physics:phyx-mini-0144:phyxminiproblems-problemphyxmini0144-lengthmagnitude, def:physics:phyx-mini-0144:phyxminiproblems-problemphyxmini0144-centimeterunitchoices}
  The nonnegative real-valued centimeter readout of a physical length
\end{definition}
\begin{proof}
  This is the declared model definition of centimeters Value.
\end{proof}
\begin{definition}[Lens Label]
  \label{def:physics:phyx-mini-0144:phyxminiproblems-problemphyxmini0144-lenslabel}
  \lean{PhyXMiniProblems.ProblemPhyXMini0144.LensLabel}
  The two lenses, in the left-to-right order shown in the figure
\end{definition}
\begin{proof}
  This is the declared model definition of Lens Label.
\end{proof}
\begin{definition}[Thin Lens Kind]
  \label{def:physics:phyx-mini-0144:phyxminiproblems-problemphyxmini0144-thinlenskind}
  \lean{PhyXMiniProblems.ProblemPhyXMini0144.ThinLensKind}
  Qualitative optical kind of a thin lens
\end{definition}
\begin{proof}
  This is the declared model definition of Thin Lens Kind.
\end{proof}
\begin{definition}[Thin Lens]
  \label{def:physics:phyx-mini-0144:phyxminiproblems-problemphyxmini0144-thinlens}
  \lean{PhyXMiniProblems.ProblemPhyXMini0144.ThinLens}
  \uses{def:physics:phyx-mini-0144:phyxminiproblems-problemphyxmini0144-lengthmagnitude, def:physics:phyx-mini-0144:phyxminiproblems-problemphyxmini0144-lenslabel, def:physics:phyx-mini-0144:phyxminiproblems-problemphyxmini0144-thinlenskind}
  A labeled thin lens with a physical focal-length magnitude
\end{definition}
\begin{proof}
  This is the declared model definition of Thin Lens.
\end{proof}
\begin{definition}[Two Lens Candle Setup]
  \label{def:physics:phyx-mini-0144:phyxminiproblems-problemphyxmini0144-twolenscandlesetup}
  \lean{PhyXMiniProblems.ProblemPhyXMini0144.TwoLensCandleSetup}
  \uses{def:physics:phyx-mini-0144:phyxminiproblems-problemphyxmini0144-lengthmagnitude, def:physics:phyx-mini-0144:phyxminiproblems-problemphyxmini0144-lenslabel, def:physics:phyx-mini-0144:phyxminiproblems-problemphyxmini0144-thinlens}
  The dimensionful data for the two-stage image formation. firstImageFromFirstLens and firstImageToSecondLens are the two positive segments into which the intermediate real image divides the lens separation. The three height fields are magnitudes, since the question asks for relative size rather than signed orientation
\end{definition}
\begin{proof}
  This is the declared model definition of Two Lens Candle Setup.
\end{proof}
\begin{definition}[stage Object Distance]
  \label{def:physics:phyx-mini-0144:phyxminiproblems-problemphyxmini0144-stageobjectdistance}
  \lean{PhyXMiniProblems.ProblemPhyXMini0144.stageObjectDistance}
  \uses{def:physics:phyx-mini-0144:phyxminiproblems-problemphyxmini0144-lengthmagnitude, def:physics:phyx-mini-0144:phyxminiproblems-problemphyxmini0144-lenslabel, def:physics:phyx-mini-0144:phyxminiproblems-problemphyxmini0144-twolenscandlesetup}
  The real-object distance used at either imaging stage
\end{definition}
\begin{proof}
  This is the declared model definition of stage Object Distance.
\end{proof}
\begin{definition}[stage Image Distance]
  \label{def:physics:phyx-mini-0144:phyxminiproblems-problemphyxmini0144-stageimagedistance}
  \lean{PhyXMiniProblems.ProblemPhyXMini0144.stageImageDistance}
  \uses{def:physics:phyx-mini-0144:phyxminiproblems-problemphyxmini0144-lengthmagnitude, def:physics:phyx-mini-0144:phyxminiproblems-problemphyxmini0144-lenslabel, def:physics:phyx-mini-0144:phyxminiproblems-problemphyxmini0144-twolenscandlesetup}
  The real-image distance used at either imaging stage
\end{definition}
\begin{proof}
  This is the declared model definition of stage Image Distance.
\end{proof}
\begin{definition}[stage Object Height]
  \label{def:physics:phyx-mini-0144:phyxminiproblems-problemphyxmini0144-stageobjectheight}
  \lean{PhyXMiniProblems.ProblemPhyXMini0144.stageObjectHeight}
  \uses{def:physics:phyx-mini-0144:phyxminiproblems-problemphyxmini0144-lengthmagnitude, def:physics:phyx-mini-0144:phyxminiproblems-problemphyxmini0144-lenslabel, def:physics:phyx-mini-0144:phyxminiproblems-problemphyxmini0144-twolenscandlesetup}
  The transverse object-height magnitude used at either imaging stage
\end{definition}
\begin{proof}
  This is the declared model definition of stage Object Height.
\end{proof}
\begin{definition}[stage Image Height]
  \label{def:physics:phyx-mini-0144:phyxminiproblems-problemphyxmini0144-stageimageheight}
  \lean{PhyXMiniProblems.ProblemPhyXMini0144.stageImageHeight}
  \uses{def:physics:phyx-mini-0144:phyxminiproblems-problemphyxmini0144-lengthmagnitude, def:physics:phyx-mini-0144:phyxminiproblems-problemphyxmini0144-lenslabel, def:physics:phyx-mini-0144:phyxminiproblems-problemphyxmini0144-twolenscandlesetup}
  The transverse image-height magnitude produced at either imaging stage
\end{definition}
\begin{proof}
  This is the declared model definition of stage Image Height.
\end{proof}
\begin{definition}[Has Depicted Lens Arrangement]
  \label{def:physics:phyx-mini-0144:phyxminiproblems-problemphyxmini0144-hasdepictedlensarrangement}
  \lean{PhyXMiniProblems.ProblemPhyXMini0144.HasDepictedLensArrangement}
  \uses{def:physics:phyx-mini-0144:phyxminiproblems-problemphyxmini0144-twolenscandlesetup}
  Both convex lenses in the primary figure are labeled as converging lenses
\end{definition}
\begin{proof}
  This is the declared model definition of Has Depicted Lens Arrangement.
\end{proof}
\begin{definition}[Has Stated Figure Readouts]
  \label{def:physics:phyx-mini-0144:phyxminiproblems-problemphyxmini0144-hasstatedfigurereadouts}
  \lean{PhyXMiniProblems.ProblemPhyXMini0144.HasStatedFigureReadouts}
  \uses{def:physics:phyx-mini-0144:phyxminiproblems-problemphyxmini0144-centimetersvalue, def:physics:phyx-mini-0144:phyxminiproblems-problemphyxmini0144-twolenscandlesetup}
  Numerical readouts printed in the primary figure: the candle is 36 cm from the first lens, the lenses are 56 cm apart, and their focal lengths are 13 cm and 16 cm. No image distance or image-size answer is assumed here
\end{definition}
\begin{proof}
  This is the declared model definition of Has Stated Figure Readouts.
\end{proof}
\begin{definition}[Has Sequential Image Geometry]
  \label{def:physics:phyx-mini-0144:phyxminiproblems-problemphyxmini0144-hassequentialimagegeometry}
  \lean{PhyXMiniProblems.ProblemPhyXMini0144.HasSequentialImageGeometry}
  \uses{def:physics:phyx-mini-0144:phyxminiproblems-problemphyxmini0144-centimetersvalue, def:physics:phyx-mini-0144:phyxminiproblems-problemphyxmini0144-twolenscandlesetup}
  Sequential-image geometry: the first real image lies between the lenses and therefore its distance from lens 1 plus its distance from lens 2 is the stated lens separation. This relation does not prescribe either segment numerically
\end{definition}
\begin{proof}
  This is the declared model definition of Has Sequential Image Geometry.
\end{proof}
\begin{definition}[Has Physical Real Image Configuration]
  \label{def:physics:phyx-mini-0144:phyxminiproblems-problemphyxmini0144-hasphysicalrealimageconfiguration}
  \lean{PhyXMiniProblems.ProblemPhyXMini0144.HasPhysicalRealImageConfiguration}
  \uses{def:physics:phyx-mini-0144:phyxminiproblems-problemphyxmini0144-centimetersvalue, def:physics:phyx-mini-0144:phyxminiproblems-problemphyxmini0144-twolenscandlesetup, def:physics:phyx-mini-0144:phyxminiproblems-problemphyxmini0144-stageobjectdistance, def:physics:phyx-mini-0144:phyxminiproblems-problemphyxmini0144-stageimagedistance}
  Nondegeneracy and real-image conditions. At each stage the positive real object distance exceeds the positive focal length, and the image and height magnitudes are nonzero
\end{definition}
\begin{proof}
  This is the declared model definition of Has Physical Real Image Configuration.
\end{proof}
\begin{definition}[Satisfies Thin Lens Equation At]
  \label{def:physics:phyx-mini-0144:phyxminiproblems-problemphyxmini0144-satisfiesthinlensequationat}
  \lean{PhyXMiniProblems.ProblemPhyXMini0144.SatisfiesThinLensEquationAt}
  \uses{def:physics:phyx-mini-0144:phyxminiproblems-problemphyxmini0144-centimetersvalue, def:physics:phyx-mini-0144:phyxminiproblems-problemphyxmini0144-lenslabel, def:physics:phyx-mini-0144:phyxminiproblems-problemphyxmini0144-twolenscandlesetup, def:physics:phyx-mini-0144:phyxminiproblems-problemphyxmini0144-stageobjectdistance, def:physics:phyx-mini-0144:phyxminiproblems-problemphyxmini0144-stageimagedistance}
  The Gaussian thin-lens equation at either lens. The division-free form f (p + q) = p q is dimensionally homogeneous and is equivalent to 1/f = 1/p + 1/q for the positive distances required above
\end{definition}
\begin{proof}
  This is the declared model definition of Satisfies Thin Lens Equation At.
\end{proof}
\begin{definition}[Satisfies Transverse Size Law At]
  \label{def:physics:phyx-mini-0144:phyxminiproblems-problemphyxmini0144-satisfiestransversesizelawat}
  \lean{PhyXMiniProblems.ProblemPhyXMini0144.SatisfiesTransverseSizeLawAt}
  \uses{def:physics:phyx-mini-0144:phyxminiproblems-problemphyxmini0144-centimetersvalue, def:physics:phyx-mini-0144:phyxminiproblems-problemphyxmini0144-lenslabel, def:physics:phyx-mini-0144:phyxminiproblems-problemphyxmini0144-twolenscandlesetup, def:physics:phyx-mini-0144:phyxminiproblems-problemphyxmini0144-stageobjectdistance, def:physics:phyx-mini-0144:phyxminiproblems-problemphyxmini0144-stageimagedistance, def:physics:phyx-mini-0144:phyxminiproblems-problemphyxmini0144-stageobjectheight, def:physics:phyx-mini-0144:phyxminiproblems-problemphyxmini0144-stageimageheight}
  The magnitude form of the transverse-magnification law at either stage: hᵢ / hₒ = q / p. It is written without division so it also preserves the physical dimensions of height times distance
\end{definition}
\begin{proof}
  This is the declared model definition of Satisfies Transverse Size Law At.
\end{proof}
\begin{definition}[final Relative Size]
  \label{def:physics:phyx-mini-0144:phyxminiproblems-problemphyxmini0144-finalrelativesize}
  \lean{PhyXMiniProblems.ProblemPhyXMini0144.finalRelativeSize}
  \uses{def:physics:phyx-mini-0144:phyxminiproblems-problemphyxmini0144-centimetersvalue, def:physics:phyx-mini-0144:phyxminiproblems-problemphyxmini0144-twolenscandlesetup}
  The dimensionless magnitude of the final image height relative to the candle
\end{definition}
\begin{proof}
  This is the declared model definition of final Relative Size.
\end{proof}
\begin{definition}[Answer Choice]
  \label{def:physics:phyx-mini-0144:phyxminiproblems-problemphyxmini0144-answerchoice}
  \lean{PhyXMiniProblems.ProblemPhyXMini0144.AnswerChoice}
  Labels of the four answer choices displayed in the source
\end{definition}
\begin{proof}
  This is the declared model definition of Answer Choice.
\end{proof}
\begin{definition}[displayed Relative Size]
  \label{def:physics:phyx-mini-0144:phyxminiproblems-problemphyxmini0144-displayedrelativesize}
  \lean{PhyXMiniProblems.ProblemPhyXMini0144.displayedRelativeSize}
  \uses{def:physics:phyx-mini-0144:phyxminiproblems-problemphyxmini0144-answerchoice}
  The dimensionless relative size printed beside each answer choice
\end{definition}
\begin{proof}
  This is the declared model definition of displayed Relative Size.
\end{proof}
\begin{definition}[recorded Answer Choice]
  \label{def:physics:phyx-mini-0144:phyxminiproblems-problemphyxmini0144-recordedanswerchoice}
  \lean{PhyXMiniProblems.ProblemPhyXMini0144.recordedAnswerChoice}
  \uses{def:physics:phyx-mini-0144:phyxminiproblems-problemphyxmini0144-answerchoice}
  Dataset metadata only: the source records answer choice C
\end{definition}
\begin{proof}
  This is the declared model definition of recorded Answer Choice.
\end{proof}
\begin{definition}[Matches Displayed Relative Size]
  \label{def:physics:phyx-mini-0144:phyxminiproblems-problemphyxmini0144-matchesdisplayedrelativesize}
  \lean{PhyXMiniProblems.ProblemPhyXMini0144.MatchesDisplayedRelativeSize}
  \uses{def:physics:phyx-mini-0144:phyxminiproblems-problemphyxmini0144-twolenscandlesetup, def:physics:phyx-mini-0144:phyxminiproblems-problemphyxmini0144-finalrelativesize, def:physics:phyx-mini-0144:phyxminiproblems-problemphyxmini0144-answerchoice, def:physics:phyx-mini-0144:phyxminiproblems-problemphyxmini0144-displayedrelativesize}
  Agreement with a displayed relative size rounded to the nearest hundredth
\end{definition}
\begin{proof}
  This is the declared model definition of Matches Displayed Relative Size.
\end{proof}
\begin{definition}[Is Nearest Displayed Choice]
  \label{def:physics:phyx-mini-0144:phyxminiproblems-problemphyxmini0144-isnearestdisplayedchoice}
  \lean{PhyXMiniProblems.ProblemPhyXMini0144.IsNearestDisplayedChoice}
  \uses{def:physics:phyx-mini-0144:phyxminiproblems-problemphyxmini0144-twolenscandlesetup, def:physics:phyx-mini-0144:phyxminiproblems-problemphyxmini0144-finalrelativesize, def:physics:phyx-mini-0144:phyxminiproblems-problemphyxmini0144-answerchoice, def:physics:phyx-mini-0144:phyxminiproblems-problemphyxmini0144-displayedrelativesize}
  A choice is correct when its displayed value is nearest to the physical ratio
\end{definition}
\begin{proof}
  This is the declared model definition of Is Nearest Displayed Choice.
\end{proof}
\begin{lemma}[stage Distances In Centimeters eq]
  \label{lem:physics:phyx-mini-0144:phyxminiproblems-problemphyxmini0144-stagedistancesincentimeters-eq}
  \lean{PhyXMiniProblems.ProblemPhyXMini0144.stageDistancesInCentimeters_eq}
  \uses{def:physics:phyx-mini-0144:phyxminiproblems-problemphyxmini0144-centimetersvalue, def:physics:phyx-mini-0144:phyxminiproblems-problemphyxmini0144-twolenscandlesetup, def:physics:phyx-mini-0144:phyxminiproblems-problemphyxmini0144-hasstatedfigurereadouts, def:physics:phyx-mini-0144:phyxminiproblems-problemphyxmini0144-hassequentialimagegeometry, def:physics:phyx-mini-0144:phyxminiproblems-problemphyxmini0144-hasphysicalrealimageconfiguration, def:physics:phyx-mini-0144:phyxminiproblems-problemphyxmini0144-satisfiesthinlensequationat}
  The figure readouts, intermediate-image geometry, and two thin-lens equations determine the successive real-image distances
\end{lemma}
\begin{proof}
  The conclusion follows from the governing relations and figure readouts named in the statement of stage Distances In Centimeters eq.
\end{proof}
\begin{lemma}[stage Relative Sizes eq]
  \label{lem:physics:phyx-mini-0144:phyxminiproblems-problemphyxmini0144-stagerelativesizes-eq}
  \lean{PhyXMiniProblems.ProblemPhyXMini0144.stageRelativeSizes_eq}
  \uses{def:physics:phyx-mini-0144:phyxminiproblems-problemphyxmini0144-centimetersvalue, def:physics:phyx-mini-0144:phyxminiproblems-problemphyxmini0144-twolenscandlesetup, def:physics:phyx-mini-0144:phyxminiproblems-problemphyxmini0144-hasstatedfigurereadouts, def:physics:phyx-mini-0144:phyxminiproblems-problemphyxmini0144-hassequentialimagegeometry, def:physics:phyx-mini-0144:phyxminiproblems-problemphyxmini0144-hasphysicalrealimageconfiguration, def:physics:phyx-mini-0144:phyxminiproblems-problemphyxmini0144-satisfiesthinlensequationat, def:physics:phyx-mini-0144:phyxminiproblems-problemphyxmini0144-satisfiestransversesizelawat}
  The magnitude laws give stage height ratios 13/23 and 92/113. Their product is the final-to-candle height ratio
\end{lemma}
\begin{proof}
  The conclusion follows from the governing relations and figure readouts named in the statement of stage Relative Sizes eq.
\end{proof}
% --- Archon declaration topology end ---
% --- Archon physics formalization source end ---
